\documentclass[aps,prd,12pt,twocolumn,superscriptaddress,floatfix,longbibliography]{revtex4-2}

\usepackage{kotex}
\usepackage{iftex}
\ifPDFTeX
  \usepackage{mathptmx}
\else
  \usepackage{fontspec}
  \setmainfont{TeX Gyre Termes}
  \setsansfont{TeX Gyre Heros}
  \setmonofont{Latin Modern Mono}
  \setmainhangulfont[
    BoldFont={UnBatangBold.ttf},
    ItalicFont={UnBatang.ttf},
    BoldItalicFont={UnBatangBold.ttf}
  ]{UnBatang.ttf}
  \setsanshangulfont[
    BoldFont={UnDotumBold.ttf},
    ItalicFont={UnDotum.ttf},
    BoldItalicFont={UnDotumBold.ttf}
  ]{UnDotum.ttf}
\fi
\usepackage{amsmath}
\usepackage{booktabs}
\usepackage{graphicx}
\usepackage{xurl}
\usepackage[a4paper, margin=0.6in]{geometry}

\setlength{\emergencystretch}{2em}
\hbadness=10000
\graphicspath{{artifacts/figures/predecoder_v2/}{artifacts/figures/predecoder/}{figure/}}
\providecommand{\href}[2]{#2}

\newcommand{\predecoderfigure}[2]{%
\IfFileExists{artifacts/figures/predecoder_v2/#1_v2.pdf}{%
    \includegraphics[width=#2]{artifacts/figures/predecoder_v2/#1_v2.pdf}%
}{%
    \IfFileExists{artifacts/figures/predecoder_v2/#1_v2.png}{%
        \includegraphics[width=#2]{artifacts/figures/predecoder_v2/#1_v2.png}%
    }{%
        \IfFileExists{artifacts/figures/predecoder/#1.pdf}{%
            \includegraphics[width=#2]{artifacts/figures/predecoder/#1.pdf}%
        }{%
            \IfFileExists{artifacts/figures/predecoder/#1.png}{%
                \includegraphics[width=#2]{artifacts/figures/predecoder/#1.png}%
            }{%
                \fbox{\parbox{0.92\linewidth}{%
                \centering Figure file missing:\\
                \texttt{artifacts/figures/predecoder_v2/#1_v2.png}}}%
            }%
        }%
    }%
}%
}

\begin{document}

\title{Design and Evaluation of a Transition-Aware Neural Pre-Decoder for Surface-Code Quantum Error Correction}
\author{Jaemin Jeon}
\affiliation{Korea University}
\date{May 2026}

\begin{abstract}
surface code 기반 양자 오류 정정에서 PyMatching은 반복 측정된 syndrome을 이용해
logical frame을 추정하는 대표적인 matching-based 기준 디코더이다. 본 연구는
PyMatching이 실패한 일부 shot에서
detector 몇 개를 국소적으로 뒤집으면 logical decision이 정답으로 돌아오는 경우가
있다는 관찰에서 출발했다. 그러나 같은 방식의 edit을 잘못 채택하면 이미 맞은 shot을
틀리게 만들 수 있으므로, 문제를 local edit 후보 생성과 후보 채택 결정으로 나누어
다루었다.

제안한 transition-aware pre-decoder는 36-channel syndrome/noise volume에서 작은
detector edit 후보를 만들고, 각 후보를 적용하기 전후의 PyMatching logical class
변화를 기준으로 후보를 점수화한다. 여기서 transition-aware는 별도의 시간 전이
모델이 아니라, local edit 전후의 downstream decoder 판단 변화를 후보 선택에
반영한다는 의미이다. 신경망은 최종 correction chain이나 logical frame을 직접
출력하지 않는다.
선택된 edit 또는 원본 syndrome만 PyMatching에 전달하고, 최종 판정은 PyMatching이
수행한다. 후보 채택 threshold는 train/validation family에서만 정했다.

학습에 사용하지 않은 \texttt{stage\_c\_corr} 평가에서 선택 모드의 pre-decoder는
d3와 d5에 대해 raw PyMatching 대비 각각 \(+0.66\) pp, \(+0.56\) pp의 평균 정확도
향상을 보였다. 다만 d5는 여덟 seed 중 두 seed에서만 보정을 채택한 보수적 결과로
해석해야 한다. d7에서는 선택 모드 향상이 \(+0.02\) pp에 그쳤다. Oracle 분석에서는
d7에도 회복 가능한 local edit 후보가 남아 있었으므로, 현재 구조의 병목은 후보 생성
자체보다 selector의 후보 정렬과 held-out noise family에 대한 일반화에 가깝다.
\end{abstract}

\keywords{quantum error correction, surface code, local edit, neural pre-decoder, PyMatching, syndrome decoding}

\maketitle

\section{서론}

Fault-tolerant quantum computation을 구현하려면, noisy stabilizer measurement로부터
물리 오류의 logical effect를 추정하는 decoder가 필요하다. surface code에서는 국소
stabilizer를 반복 측정하고, 시간에 따라 변하는 syndrome을 detector event로
표현한다~\cite{surfacecode}. 이 detector event는 실제 error path를 직접 알려
주지는 않지만, 관측된 syndrome과 양립 가능한 error chain의 후보를 제한한다.

PyMatching을 포함한 matching-based decoder는 이러한 inference 문제에서 안정적인
기준 성능을 제공한다~\cite{pymatching}. 실제 실험을 진행하면서 먼저 확인한 점은,
PyMatching이 틀린 일부 샘플이 detector 몇 개를 국소적으로 뒤집는 것만으로도
정답 logical frame으로 돌아올 수 있다는 사실이었다. 반대로 같은 방식의 보정을
잘못 적용하면 이미 맞은 샘플을 틀리게 만들 수도 있었다. 이 관찰에 따라
신경망이 최종 logical class를 직접 예측하도록 만들기보다, PyMatching에 들어가기
전 syndrome을 아주 제한적으로 수정할지를 판단하는 문제로 범위를 좁혔다.

이 설정은 neural pre-decoding 계열의 문제 설정과 가깝다. Pre-decoder는 완전한
decoder 역할을 맡기보다, 뒤따르는 global decoder에 들어가기 전 syndrome을 줄이거나
수정하는 앞단 모듈로 동작한다. NVIDIA의 AI-based pre-decoder 연구는 이 방향을 대규모
surface-code setting에서 보여준 선행 연구이다~\cite{nvidiaising}. 여기서는
대규모 하드웨어 실험과의 직접 비교보다는, 작은 거리에서 local edit 후보를 만들고
그 후보를 채택하거나 버리는 결정이 PyMatching 성능에 어떤 영향을 주는지 분해했다.

구현상으로는 3D residual syndrome trunk를 먼저 학습하고, PyMatching 실패 샘플에서
찾은 local edit motif를 후보로 만들었다. 이후 \texttt{CandidateEditSelector}가
각 후보를 점수화하도록 학습했는데, 이때 label은 detector mask가 맞는지가 아니라
그 후보를 적용한 뒤 PyMatching의 \texttt{logical\_class4} 판단이 개선되는지로
정했다. 결과 해석에서도 raw PyMatching, 후보를 그대로 적용한 결과, 안전 정책을
거친 결과, 그리고 후보 공간 안에서 가능한 oracle 결과를 분리했다. 이렇게 나누어
보아야 d3/d5의 실제 향상과 d7에서의 실패 원인을 같은 평가 축에서 비교할 수 있다.

\section{배경 및 관련 연구}

\subsection{surface code와 matching decoder}

surface code는 2차원 격자 위의 data qubit와 측정 qubit를 이용하여 stabilizer
measurement를 반복한다. 같은 stabilizer 위치에서 측정값이 시간에 따라 바뀌면
detector event가 발생하며, boundary 및 final measurement에 의해 닫히는 event도
decoder graph에 포함된다. Decoder는 이 event pattern과 code boundary를 함께
사용하여 가능한 error chain을 추정하고, 그 chain이 유도하는 logical frame을
판단한다.

PyMatching은 이 문제를 sparse matching graph 위의 minimum-weight matching 문제로
정식화한다~\cite{pymatching}. 여기서 PyMatching은 세 가지 역할을 갖는다.
첫째, 보정을 적용하지 않은 원본 syndrome에 대한 기준 디코더이다. 둘째, 후보
edit의 benefit/harm label을 산정하는 평가 기준이다. 셋째, neural pre-decoder가
선택한 syndrome을 최종적으로 판정하는 디코더이다. 본 논문의 성능 비교는 신경망
보정기와 PyMatching을 결합한 전체 pipeline의 결과이다.
본문의 \texttt{logical\_class4}는 PyMatching의 logical 판단을 네 개 class로 정리한
평가 label을 뜻한다. 문맥에 따라 logical frame 또는 logical decision이라고도
부르지만, 모두 같은 최종 logical 판단을 가리킨다.

\subsection{신경망 디코딩과 사전 디코딩}

신경망을 이용한 surface code 디코딩은 크게 두 계열로 나눌 수 있다. 첫 번째는
syndrome에서 logical class 또는 correction field를 직접 예측하는 replacement-style
neural decoder이다. AlphaQubit과 AlphaQubit 2는 이 방향의 대표적인 대규모
연구이다~\cite{alphaqubit,alphaqubit2}. 또한 convolutional neural decoder는
surface-code lattice의 local pattern을 직접 활용하는 direct decoder 계열의
예시이다~\cite{convnnsurface}. 두 번째는 기존 global decoder를 유지하되, 그
앞단에서 syndrome을 보정하거나 residual problem을 줄이는 pre-decoding 방식이다.
NVIDIA Ising-Decoding은 neural pre-decoder와 뒤따르는 decoder를 결합한 구조적
선행 연구로 볼 수 있다~\cite{nvidiaising}.

본 논문은 두 번째 계열에 속한다. 직접 neural classifier를 사용해 final
logical class를 예측하는 방식은 작은 거리에서는 시도 가능하지만, 거리 증가와
noise-family shift에 대해 안정적인 기준을 제공하기 어렵다. 반면 pre-decoder
구조는 PyMatching의 안정성을 유지하면서, 특정 local syndrome pattern에 대한
입력 수정만 제한적으로 허용한다. 여기서는 neural network의 역할을
``local edit 후보 선택''으로 둔다.

관련 연구와의 비교에서 중요한 구분은 ``무엇을 신경망이 직접 결정하는가''이다.
AlphaQubit 계열과 CNN decoder는 syndrome으로부터 final decision 또는 correction을
직접 예측하는 방향에 가깝다~\cite{alphaqubit,alphaqubit2,convnnsurface}. 반면 본
연구에서는 최종 matching 문제를 PyMatching에 남겨 두고, 신경망은 그 입력을
국소적으로 수정할지 결정한다. NVIDIA Ising-Decoding은 pre-decoder와 후속 decoder를
결합한다는 점에서 가장 가까운 구조적 참고점이지만~\cite{nvidiaising}, 여기서는
작은 거리에서 후보 생성과 후보 선택의 실패 경로를 더 세밀하게 본다. Google Willow
결과는 hardware-scale QEC의 문맥을 제공하지만~\cite{willowthreshold}, 본 논문의
수치와 직접 비교하는 대상은 아니다.

\begin{table*}[t]
\caption{관련 연구와 본 논문의 위치.}
\label{tab:related-work}
\centering
\scriptsize
\begin{tabular*}{\textwidth}{@{\extracolsep{\fill}}llll}
\toprule
work & learned decision & final decoder & relation \\
\midrule
AlphaQubit family~\cite{alphaqubit,alphaqubit2} & logical/correction decision & neural decoder & direct final-decision model \\
CNN surface decoder~\cite{convnnsurface} & local-pattern correction & CNN decoder & direct decoder baseline \\
NVIDIA Ising-Decoding~\cite{nvidiaising} & pre-decoding decision & residual decoder & pre-decoder before residual decoder \\
Near-term neural decoder~\cite{neartermnn} & experimental syndrome mapping & neural decoder & experimental syndrome-to-logical mapping \\
Google Willow~\cite{willowthreshold} & hardware-scale QEC evidence & hardware decoder stack & hardware-scale motivation only \\
This work & local detector edit & PyMatching & local-edit selection before PyMatching \\
\bottomrule
\end{tabular*}
\end{table*}

\section{데이터셋 및 노이즈 환경}

실험은 Stim 기반으로 생성한 rotated surface-code memory circuit과 자체 noise
family를 사용해 수행했다~\cite{stim}. Noise family는 단계적으로 구성했다.
\texttt{stage\_a\_si1000}은 surface-code resource estimate에서 사용된
superconducting-inspired 설정을 따른 기본 SI1000 계열 noise 환경이고
~\cite{honeycombmemory},
\texttt{stage\_b\_local}은 local heterogeneity를 추가한 환경이며,
\texttt{stage\_c\_corr}는 correlated stray surrogate를 포함한 held-out 평가
환경이다. 본 논문의 정량 주장은 \texttt{stage\_c\_corr}에서의 성능을 기준으로
한다.

모델 입력은 detector event만으로 구성하지 않았다. Detector mask, geometry
metadata, noise-family indicator, distance/round statistics, event fraction,
physical-noise summary 등을 함께 포함하여 36-channel syndrome/noise volume을
구성했다. 채널 의미는 code distance와 무관하게 고정하고, spatial/temporal
volume의 크기만 distance에 따라 달라진다. 예를 들어 d3, d5, d7의 입력 shape은
각각 \([36,4,4,4]\), \([36,6,6,6]\), \([36,8,8,8]\)이다.

각 distance의 원본 manifest는 family별 \(2048\) shots로 생성했다. 학습용 local edit
target은 train/validation family의 고정된 \(1024\)-shot subset에서 생성했다.
본문에 보고한 주요 결과는 학습에 사용하지 않은 \texttt{stage\_c\_corr}의
\(1024\) shots/seed에서 계산했고, held-out local edit 탐색은 평가 후 oracle 회복
가능성 분석에만 사용했다. 각 train family에서는 고정된 validation split
\(154\) shots를 사용했고, 나머지 샘플이 trunk와 selector 학습에 사용된다.

\begin{table*}[t]
\caption{Distance별 데이터 구성. 원본 dataset은 family별 \(2048\) shots이고, 학습용
local edit target은 train/validation family의 고정된 \(1024\)-shot subset에서
생성했다. Held-out family의 \(1024\) shots는 최종 평가와 oracle 분석에만 사용했다.}
\label{tab:dataset-config}
\centering
\scriptsize
\begin{tabular*}{\textwidth}{@{\extracolsep{\fill}}lccccc}
\toprule
distance & rounds & input shape & train/validation families & held-out family & reported seeds \\
\midrule
d3 & 3 & \([36,4,4,4]\) & A/B, \(1024\) shots each & C, \(1024\) shots & \(0..7\) \\
d5 & 5 & \([36,6,6,6]\) & A/B, \(1024\) shots each & C, \(1024\) shots & \(0..7\) \\
d7 & 7 & \([36,8,8,8]\) & A/B, \(1024\) shots each & C, \(1024\) shots & \(0..57\) \\
\bottomrule
\end{tabular*}
\end{table*}

\begin{table*}[t]
\caption{실험에 사용한 noise family.}
\label{tab:noise-family}
\centering
\footnotesize
\begin{tabular*}{\textwidth}{@{\extracolsep{\fill}}llll}
\toprule
family & noise structure & role & reported use \\
\midrule
\texttt{stage\_a\_si1000} & SI1000-style baseline physical noise & baseline family & train/validation \\
\texttt{stage\_b\_local} & Stage A plus local heterogeneity & structured-noise family & train/validation \\
\texttt{stage\_c\_corr} & Stage B plus correlated stray surrogate & primary test family & held-out evaluation \\
\bottomrule
\end{tabular*}
\end{table*}

\section{제안 방법}

제안 방법은 PyMatching 앞단에 배치되는 신경망 사전 디코더이다.
전체 추론 흐름은 Figure~\ref{fig:pipeline}과 같다. 신경망은 detector-event
volume에서 small local edit 후보를 만들고, 각 후보가 PyMatching 결과에
유익한지 또는 해로운지 평가한다. 안전 조건을 통과한 후보만 syndrome에
적용되며, 그렇지 않은 경우 raw no-edit syndrome이 그대로 PyMatching에 전달된다.
Selector는 raw no-edit syndrome을 넣었을 때의 PyMatching logical class와, 특정
local edit을 적용한 뒤의 PyMatching logical class 사이의 변화를 학습 신호와
candidate feature로 사용한다. 즉 detector를 복원하는 문제보다, 그 edit이 최종
logical 판단을 어느 방향으로 바꾸는지를 기준으로 후보를 고른다.

\begin{figure*}[t]
\centering
\predecoderfigure{fig1_predecoder_pipeline}{\textwidth}
\caption{제안하는 neural pre-decoder와 PyMatching 결합 구조. 신경망은 local
syndrome edit의 채택 여부만 결정하고, 최종 logical frame은 PyMatching이
예측한다.}
\label{fig:pipeline}
\end{figure*}

\subsection{SyndromeEditPreDecoder}

\texttt{SyndromeEditPreDecoder}는 3D convolution stem과 세 개의 residual 3D
convolution block으로 구성된다. 최종 d3/d5 설정에서는 hidden channel
\(24\), dense hidden dimension \(64\), dropout \(0.1\)을 사용했으며, 해당 설정의
parameter 수는 \(118834\)이다. Trunk는 detector-level \texttt{edit\_logits},
shot-level \texttt{needs\_edit\_logits}, 그리고 pooled shot feature를 출력한다.
이 중 pooled feature와 local candidate feature가 patch-head selector의 입력으로
사용된다.

Trunk 학습은 Adam optimizer를 사용했다. 기본 학습 설정은 batch size \(64\),
learning rate \(10^{-3}\), 최대 \(8\) epochs이다. Loss는 세 부분으로 구성했다.
첫째, known local-edit target이 있는 detector 위치에 대해 masked binary
cross-entropy를 적용했다. 둘째, shot 단위로 edit이 필요한지를 예측하는
\texttt{needs\_edit\_logits}에 binary cross-entropy를 적용했다. 셋째, edit이
필요한 샘플에서 target edit이 identity 후보보다 높은 점수를 갖도록 margin \(1.0\)
의 ranking loss를 더했다. 실제 보고값에는 validation-family metric
이 가장 좋은 epoch의 checkpoint를 사용했다.

\subsection{Local motif candidate}

모델은 많은 detector를 한 번에 뒤집는 대규모 syndrome 보정을 허용하지 않는다.
Target 생성 단계에서는 raw PyMatching이 틀린 샘플에 대해서만 국소 탐색을 수행했다.
탐색 반경은 시간 방향 \(1\), 공간 방향 \(1\)로 두었고, 한 후보가 뒤집을 수 있는
detector 수는 최대 \(2\)개로 제한했다. 각 incorrect shot에서 최대 \(20\)개의 local
variant를 검사하여, PyMatching 결과를 정답 \texttt{logical\_class4}로 되돌리는
edit을 target으로 기록했다. 예를 들어 held-out \texttt{stage\_c\_corr}에서는 d3의
raw 실패 73개 중 65개, d5의 raw 실패 114개 중 92개, d7의 raw 실패 130개 중
114개가 이런 local edit 탐색으로 oracle 수준에서는 회복 가능했다.
이 held-out family의 local edit 탐색 결과는 학습, 후보 motif 선택, threshold
결정에 사용하지 않았고, 후보 공간의 oracle 회복 가능성을 평가하기 위한 사후
분석으로만 사용했다.

최종 후보 구성에서는 train/validation family에서 얻은 target edit들을 그대로 모두
쓰지 않고 \texttt{local\_motif\_selector}로 후보를 줄였다. 관측된 local edit
pattern을 motif로 모은 뒤 최대 \(16\)개 class만 유지하고, inference 때는 trunk의 detector별 edit
확률 합이 큰 motif placement를 top-\(32\)개까지 후보로 넣었다. Identity, 즉 아무
보정도 하지 않는 후보는 항상 포함된다. 이 구조 때문에 raw PyMatching은 외부
baseline일 뿐 아니라, 선택 모드 내부에서 안전하게 되돌아갈 수 있는 기준 경로가
된다.

\subsection{Benefit/harm selector}

\texttt{CandidateEditSelector}는 각 local edit 후보에 scalar score를 부여한다.
후보의 label은 해당 edit을 적용한 뒤 PyMatching의 \texttt{logical\_class4} 판단이
raw no-edit보다 개선되었는지 여부로 정의한다.
Figure~\ref{fig:architecture}는 이 구조를 요약한다.

Selector 입력은 shot-level pooled feature와 candidate feature의 결합이다. Candidate
feature에는 edit weight, detector별 edit probability의 합/평균/최댓값, motif
빈도, detector 좌표 요약, 주변 event/evidence, 그리고 \(3\times3\times3\) local
patch에서 읽은 event/probability 값이 포함된다. Patch-head는 이 local patch slice를
작은 MLP로 먼저 embedding한 뒤, 나머지 candidate feature 및 shot feature와
concatenate한다. Selector 학습은 Adam optimizer, learning rate \(5\times10^{-4}\),
최대 \(6\) epochs로 수행했다. 기본 loss는 후보 group 안에서 가장 좋은 후보를
고르는 cross-entropy이고, 여기에 유익한 nonzero edit이 해로운 nonzero edit보다
높은 점수를 갖도록 하는 pairwise loss를 weight \(1.0\), margin \(0.5\)로
추가했다. 이미 맞은 shot을 망치는 후보에는 harm weight \(2.0\), raw가
틀린 shot에서 여전히 틀린 nonzero 후보에는 miss weight \(0.25\)를 부여했다.

\begin{figure*}[t]
\centering
\predecoderfigure{fig2_model_architecture}{\textwidth}
\caption{Neural pre-decoder의 구조. \texttt{SyndromeEditPreDecoder}가 shared 3D
feature volume을 만들고 edit, needs-edit, pooled-shot head를 출력한다. 이후
patch-head \texttt{CandidateEditSelector}가 candidate feature를 점수화하고,
선택된 보정 또는 원본 유지 경로가 PyMatching으로 전달된다.}
\label{fig:architecture}
\end{figure*}

\subsection{Selected-mode safety}

Candidate branch가 항상 최종 결과에 반영되는 것은 아니다. 채택 조건과 threshold는
train/validation family에서만 정했고, held-out \texttt{stage\_c\_corr} 결과는
threshold나 seed별 채택 여부를 정하는 데 사용하지 않았다. Validation family에서
후보 적용이 raw no-edit 대비 충분한 validation gain을 보이지 않거나, 해로운 후보
수와 selector margin 조건을 넘어서면 선택 모드는 보정을 적용하지 않고 원본
syndrome을 그대로 PyMatching에 전달한다. 이 후보 우선 안전 정책은 local edit이
일부 seed에서 PyMatching을 악화시킬 수 있다는 관찰에 기반한다.

구체적으로 validation family에서 candidate branch가 raw no-edit보다 충분히 좋아야
한다. Validation delta가 \(+0.02\) 이상이면 충분한 채택 근거로 보고 그대로
받아들였다. 약한 양의 증거는 delta \(+0.005\) 이상, selector margin \(0.5\) 이상,
harmed count가 \(2\) 이하일 때만 채택했다. 그 밖의 경우에는 보정을 적용하지 않는다.
d5에서 seed 2와
seed 3만 local edit을 채택하고 나머지 seed가 원본 syndrome을 유지한 것은 이
조건을 그대로 적용한 결과이다.

\section{실험 설정}

비교 대상은 네 가지 경로로 정의했다. Raw PyMatching은 원본 syndrome을 그대로
사용한 기준 결과이다. Selected predecoder는 안전 정책이 선택한 local edit 또는
수정하지 않은 syndrome을 적용한 제안 방법의 최종 결과이다. Candidate branch는
selector가 고른 후보를 안전 정책 없이 바로 적용했을 때의 성능이며, target
local-edit oracle은 정해진 후보 공간 안에서 가능한 최선의 local edit 성능이다.

주요 metric은 학습에 사용하지 않은 \texttt{stage\_c\_corr}에서의
\texttt{logical\_class4} accuracy이다. Candidate branch와 selected branch를
분리해 보고한 이유는, 후보 공간의 잠재력과 실제 채택 정책의 안정성을 구분하기
위해서이다. 이 구분이 없으면 d7에서 관측되는 oracle 기준 회복 가능성과 낮은
선택 모드 향상의 차이를 설명하기 어렵다.

\section{실험 결과}

\subsection{Main accuracy}

Table~\ref{tab:main-results}와 Figure~\ref{fig:main-accuracy}는 d3, d5, d7의
주요 결과를 보여준다. d3와 d5에서는 선택 모드의 pre-decoder가 raw PyMatching보다
높은 평균 accuracy를 보였다. 반면 d7에서는 정확도 향상이 \(+0.02\) pp에 그쳐,
현재 학습된 selector가 의미 있는 회복을 만들었다고 보기는 어렵다.

\begin{table*}[t]
\caption{Held-out \texttt{stage\_c\_corr}에서의 main accuracy.}
\label{tab:main-results}
\centering
\scriptsize
\begin{tabular*}{\textwidth}{@{\extracolsep{\fill}}lcccccc}
\toprule
distance & seeds & raw & selected & candidate & oracle & gain \\
\midrule
d3 & \(0..7\) & 92.87\% & 93.53\% & 93.53\% & 99.22\% & \(+0.66\) pp \\
d5 & \(0..7\) & 88.87\% & 89.43\% & 89.18\% & 97.85\% & \(+0.56\) pp \\
d7 & \(0..57\) & 87.30\% & 87.32\% & 87.15\% & 98.44\% & \(+0.02\) pp \\
\bottomrule
\end{tabular*}
\end{table*}

\begin{figure*}[t]
\centering
\predecoderfigure{fig3_main_accuracy_comparison}{\textwidth}
\caption{거리별 held-out accuracy. d3/d5에서는 선택 모드의 pre-decoder가 raw
PyMatching을 넘지만, d7에서는 oracle 기준 회복 가능성이 큰데도 실제 선택 결과의
회복은 거의 나타나지 않는다.}
\label{fig:main-accuracy}
\end{figure*}

\subsection{d3/d5 robustness}

d3는 seed \(0..7\) 전부에서 양의 정확도 변화를 보였다. 아래 표의 seed classes는
selected delta가 positive/neutral/negative인 seed 수를 뜻한다. d3에서도 shot
단위로는 개선된 경우와 악화된 경우가 함께 존재하지만, 여덟 seed의 net effect는
모두 양수였다. d5는 두 seed에서만 local selector를 채택하고, 나머지 여섯 seed에서는
원본 syndrome을 유지했다. d5의 평균 향상은 유익한 경우에만 보정을 채택한
selected-mode 결과로 제한된다.

\begin{table}[t]
\caption{d3/d5 seed \(0..7\) robustness summary. Seed classes are
positive/neutral/negative seed counts.}
\label{tab:d3d5-robustness}
\centering
\footnotesize
\begin{tabular*}{\columnwidth}{@{\extracolsep{\fill}}lccc}
\toprule
distance & seed classes & gain & improved/harmed \\
\midrule
d3 & \(8/0/0\) & \(+0.66\) pp & \(205/151\) \\
d5 & \(2/6/0\) & \(+0.56\) pp & \(56/10\) \\
\bottomrule
\end{tabular*}
\end{table}

Seed-level bootstrap 95\% confidence interval은 d3가 약 \([+0.45,+0.85]\) pp,
d5가 약 \([0.00,+1.37]\) pp이다. Exact sign test에서도 d3의 one-sided p-value는
\(0.0039\)인 반면, d5는 \(0.25\)이다. 따라서 d3는 검사한 seed 전부에서 양의
결과를 보인 사례로 제시할 수 있지만, d5는 positive mean과 선택 조건의 안정성을
중심으로 해석해야 한다.

\begin{table}[t]
\caption{d3/d5 paired seed-level statistics.}
\label{tab:paired-stats}
\centering
\footnotesize
\begin{tabular*}{\columnwidth}{@{\extracolsep{\fill}}lccc}
\toprule
distance & seed classes & sign p & sign-flip p \\
\midrule
d3 & \(8/0/0\) & 0.0039 & 0.0039 \\
d5 & \(2/6/0\) & 0.2500 & 0.2500 \\
\bottomrule
\end{tabular*}
\end{table}

\begin{table*}[t]
\caption{d5 seed별 선택 결과와 원본 유지 여부.}
\label{tab:d5-fallback}
\centering
\scriptsize
\begin{tabular*}{\textwidth}{@{\extracolsep{\fill}}clrrl}
\toprule
seed & 선택 경로 & selected delta & candidate delta & 해석 \\
\midrule
0 & 원본 유지 & \(+0.00\) pp & \(+0.00\) pp & 보정 미채택 \\
1 & 원본 유지 & \(+0.00\) pp & \(+0.00\) pp & 보정 미채택 \\
2 & 보정 채택 & \(+2.15\) pp & \(+2.15\) pp & 유익한 보정 \\
3 & 보정 채택 & \(+2.34\) pp & \(+2.34\) pp & 유익한 보정 \\
4 & 원본 유지 & \(+0.00\) pp & \(-0.78\) pp & 해로운 후보 차단 \\
5 & 원본 유지 & \(+0.00\) pp & \(+0.00\) pp & 보정 미채택 \\
6 & 원본 유지 & \(+0.00\) pp & \(-1.17\) pp & 해로운 후보 차단 \\
7 & 원본 유지 & \(+0.00\) pp & \(+0.00\) pp & 보정 미채택 \\
\bottomrule
\end{tabular*}
\end{table*}

Table~\ref{tab:d5-fallback}는 d5의 평균 향상이 seed 2와 seed 3에서 발생했음을
보여준다. 반대로 seed 4와 seed 6의 candidate branch는 해로운 후보였지만,
선택 모드는 보정을 적용하지 않음으로써 해당 성능 하락을 막았다.

\subsection{Candidate branch와 selected branch}

Candidate branch는 후보 공간의 잠재력을 보여주지만, 안전 정책 없이 적용하면
해로운 보정도 함께 반영된다. 특히 d7에서는 candidate branch의 improved/harmed
shot 수가 \(161/251\)로 harmed 쪽이 더 많다. 선택 모드는 대부분 보정을
적용하지 않아 큰 성능 하락을 막지만, 동시에 oracle 기준의 회복 가능성을 충분히
활용하지도 못한다. 아래 결과에서는 candidate, selected, oracle branch를 구분해
표기한다.
d3에서는 candidate와 selected accuracy가 모두 93.53\%로 같아, fallback이 결과를
바꾸지 않았다. 반면 selected와 candidate의 차이는 d5에서 \(+0.25\) pp, d7에서
\(+0.17\) pp이며, 이는 safety policy가 해로운 후보를 차단해 얻은 효과를 보여준다.

\section{d7 한계 분석}

d7의 선택 모드 정확도 변화는 \(+0.02\) pp에 불과하다. 그러나 target local-edit
oracle accuracy는 98.44\%이고, 모든 58개 seed에 oracle 기준의 회복 가능성이
존재한다. 따라서 d7의 한계를 후보 공간 안에 유익한 local edit이 없어서 생긴
문제로 해석하는 것은 적절하지 않다.

더 직접적인 문제는 selector가 학습에 사용하지 않은 평가 환경에서 유익한 후보와
해로운 후보를 안정적으로 구분하지 못한다는 점이다. d7 candidate outcome은
positive \(6\), neutral \(35\), harmful \(17\)로 나뉘며, 후보를 안전 정책 없이
적용하면 평균적으로 \(-0.15\) pp 하락했다. Figure~\ref{fig:d7-oracle}은 oracle
기준 회복 가능성과 false-positive 구조를 함께 보여준다.

\begin{figure*}[t]
\centering
\predecoderfigure{fig4_d7_oracle_gap_false_positive}{\textwidth}
\caption{d7의 oracle 기준 회복 가능성과 false-positive 후보 문제. 후보 공간에는
유익한 edit이 남아 있지만, learned selector가 validation-positive branch를
held-out에서 안정적으로 고르지 못한다.}
\label{fig:d7-oracle}
\end{figure*}

Validation 결과와 held-out 결과의 불일치는 selector 일반화 문제가 주요 병목임을
가리킨다.
Validation-positive candidate seed 22개 중 held-out에서 해로운 결과로 바뀐 seed는
13개이고, positive로 남은 seed는 5개뿐이다. 즉 validation evidence만으로 후보
채택을 결정하면 false-positive에 해당하는 해로운 edit을 충분히 차단하기 어렵다.
여기서 false-positive는 validation-positive였지만 held-out에서 harmful로 바뀐
seed를 뜻하며, held-out neutral 4개는 harmful false-positive 비율의 분자에 넣지
않았다. Validation-positive seed 중 harmful held-out outcome 비율은
\(13/22=59.09\%\)이다.

\begin{table}[t]
\caption{d7 validation class와 held-out outcome의 불일치.}
\label{tab:d7-mismatch}
\centering
\footnotesize
\begin{tabular*}{\columnwidth}{@{\extracolsep{\fill}}lccc}
\toprule
validation class & harmful & neutral & positive \\
\midrule
neutral & 4 & 31 & 1 \\
positive & 13 & 4 & 5 \\
\bottomrule
\end{tabular*}
\end{table}

\begin{figure*}[t]
\centering
\predecoderfigure{fig5_d7_validation_heldout_scatter}{\textwidth}
\caption{d7 validation delta와 held-out candidate delta의 seed-level scatter.
Validation-positive evidence가 held-out correlated noise에서 그대로 유지되지
않는다는 점이 핵심이다.}
\label{fig:d7-scatter}
\end{figure*}

해로운 candidate seed는 총 17개였고, 선택 모드의 guard는 이 \(17/17\)개에서
모두 보정을 적용하지 않고 원본 syndrome을 유지했다. 즉 안전 정책은 d7에서
성능을 크게 올리지는 못했지만, 해로운 후보가 실제 출력으로 전달되는 것은 막았다.

d7에 대해서는 threshold sweep도 추가로 수행했지만, 현재 model family 안에서는
합리적인 해결책을 찾지 못했다. 약 \(18.3\)만 개의 scalar adoption-threshold
policy를 검사했지만, 사전에 정의한 preserve/recover/block seed-level 조건을
동시에 만족한 policy는 없었다. 정확한 sentinel seed 집합은 재현성 패키지에
기록했다.
유익한 후보와 해로운 후보를 직접 구분하도록 학습 목표를 바꾸거나, 후보 간
호환성을 반영한 pairwise top-k를 적용한 경우에도 true positive 보존과 false
positive 차단을 동시에 만족하지 못했다. 따라서 d7의 다음 개선은 단순 threshold
조정이 아니라 selector의 후보 정렬과 일반화 목표를 다시 설계하는 방향이어야 한다.

\begin{figure*}[t]
\centering
\predecoderfigure{fig6_oracle_recovery_distribution}{\textwidth}
\caption{Seed-level oracle recovery. d3/d5는 oracle 기준 회복 가능성의 일부를 실제
selected 결과로 회복하지만, d7 선택 모드는 candidate-oracle 기준의 회복 가능성을
거의 활용하지 못한다.}
\label{fig:oracle-recovery}
\end{figure*}

\section{논의}

d3/d5 결과에서 중요한 점은 보정 자체보다 채택 조건이다. d5에서는 candidate
branch만 보면 해로운 seed가 있었지만, validation에서 정한 조건을 통과한 seed 2와
seed 3에서만 보정을 적용했다. 이 결과는 제한된 후보 집합에서 PyMatching에 도움이
되는 경우를 selector가 가려낼 수 있었다는 정도로 읽어야 한다.

d7에서는 같은 절차가 거의 이득을 만들지 못했다. Oracle 분석에서는 회복 가능한
local edit이 남아 있었지만, validation-positive 후보 중 상당수가 held-out에서
해로운 결과로 바뀌었다. 즉 후보 생성보다 selector의 noise-family 일반화가 더
직접적인 병목이었다.

본 논문의 정량 평가는 Stage A/B/C noise family 안에서 해석해야 한다. 이 범위는
pre-decoder의 후보 생성과 후보 채택을 분리해 보기 위해 정한 실험 범위이며,
대규모 hardware decoder 성능을 주장하는 범위는 아니다.

\section{결론}

본 논문에서는 PyMatching을 대체하지 않고, 그 앞단에서 작은 detector edit을
선택하는 신경망 사전 디코더를 구현했다. d3/d5의 held-out
\texttt{stage\_c\_corr}에서는 각각 \(+0.66\) pp, \(+0.56\) pp의 평균 향상을
얻었다. 앞의 seed-level 분석에서 보듯 d5 향상은 두 seed만 보정을 채택한
selected-mode 개선으로 남는다.

d7에서는 \(+0.02\) pp에 그쳤다. 후보 공간에는 회복 가능한 edit이 있었지만,
validation에서 좋아 보인 후보가 held-out에서 해로운 경우가 많았다. 따라서 다음
개선은 더 큰 후보 집합보다 selector target과 validation-to-held-out 일반화 문제를
다루는 쪽이 되어야 한다.

\section*{재현성 관리}

본 연구의 표와 그림은 동일한 평가 산출물에서 생성하였다. 최종 결과표, d3/d5
seed별 비교 통계, oracle 기준 회복률 분포, d7의 해로운 후보 분류를 각각 별도
기록으로 보존했고, 서로 다른 표와 그림이 같은 수치에 기반하는지 consistency
check로 점검했다. 논문에 사용한 최종 산출물은 \(37/37\)개의 consistency check를
통과했으며, 실패 항목은 \(0\)개였다. 관련 summary와 재현성 메모는
\texttt{artifacts/eval/nn/} 및 \texttt{PREDECODER\_REPRODUCIBILITY\_PACKAGE.md}에
보존했다. Hyperparameter sensitivity 기록은 본문 표에는 포함하지 않고
\texttt{PREDECODER\_HYPERPARAMETER\_SENSITIVITY.md}에 별도로 보존했다.

\bibliographystyle{unsrtnat}
\bibliography{main}

\end{document}
